% Figure from MATH 590 notes, section 15. Compile with the notes preamble.
\begin{tikzpicture}[x=2.3cm, y=0.36cm]
  % N = { |x| < 1/(|y|+1) }
  \fill[figB!12] plot[domain=-6.5:6.5, samples=120, variable=\t] ({1/(abs(\t)+1)},\t)
    -- plot[domain=6.5:-6.5, samples=120, variable=\t] ({-1/(abs(\t)+1)},\t) -- cycle;
  \draw[figB, thick, dashed] plot[domain=-6.5:6.5, samples=120, variable=\t] ({1/(abs(\t)+1)},\t);
  \draw[figB, thick, dashed] plot[domain=-6.5:6.5, samples=120, variable=\t] ({-1/(abs(\t)+1)},\t);
  \draw[-stealth, thin, black!70] (-1.3,0) -- (1.35,0) node[right, font=\scriptsize] {$x$};
  \draw[-stealth, thin, black!70] (0,-6.6) -- (0,6.9) node[above, font=\scriptsize] {$y$};
  % the slice {0} x R
  \draw[figB, very thick] (0,-6.5) -- (0,6.5);
  % a would-be tube (-eps, eps) x R with eps = 1/4, sticking out for |y| > 3
  \draw[figR, thick, dashed] (-0.25,-6.5) -- (-0.25,6.5) (0.25,-6.5) -- (0.25,6.5);
  \foreach \s in {1,-1} {
    \fill[figR!35] plot[domain=3:6.5, samples=40, variable=\t] ({1/(\t+1)},\s*\t) -- (0.25,\s*6.5) -- (0.25,\s*3) -- cycle;
    \fill[figR!35] plot[domain=3:6.5, samples=40, variable=\t] ({-1/(\t+1)},\s*\t) -- (-0.25,\s*6.5) -- (-0.25,\s*3) -- cycle;
  }
  \node[font=\scriptsize, black!70, anchor=north west, inner sep=1pt] at (0.26,-0.1) {$\varepsilon$};
  \node[font=\scriptsize, black!70, anchor=north east, inner sep=1pt] at (-0.26,-0.1) {$-\varepsilon$};
  \node[font=\scriptsize, black!70] at (1.0,-0.55) {$1$};
  \node[font=\small, figB] at (0.72,1.6) {$N$};
  \node[font=\scriptsize, figB, align=left, right] at (0.45,4.6) {$|x|=\dfrac{1}{|y|+1}$};
  \node[font=\scriptsize, figR, right, align=left] at (0.3,-4.9) {tube leaves $N$\\ once $|y|>\tfrac1\varepsilon-1$};
\end{tikzpicture}
